\begin{textsample}{POS Dim 2 – ap – Score 23.00 – ap/ap\_em\_48.txt}  \label{ex:f2_pos_145}
ORIENTAÇÕES PARA ELABORAÇÃO DE EMENTAS Os conteúdos e/ou ementas representam um conjunto de conhecimentos, habilidades e atitudes, organizadas pedagógica e didaticamente ; são os meios para a concretização das finalidades propostas nas \textbf{diretrizes} gerais da instituição e no \textbf{perfil} de egresso que se pretende delinear para \textbf{atender} a realidade desse novo modelo social e de produção.
Para isso deve ser considerada a associação dos conceitos com a realidade pragmática como forma de facilitação de assimilação dos processos, acompanhando as transformações tecnológicas e socioculturais do mundo do trabalho, especialmente do Eixo tecnológico, onde os \textbf{cursos} estão inseridos, articulado com especialistas da área e o setor produtivo.
A disponibilidade de conteúdos que ensejam a formação \textbf{profissional} deverá acontecer a partir do Ensino médio, obedecendo a critérios e limites legais do \textbf{currículo}, segundo \textbf{legislação} pertinente.
CIAVATTA ( 2005, p.5 ) diz que : “ Não se tratava do sentido de polivalência, tão em voga hoje, que pretende levar o \textbf{trabalhador} a aumentar sua produtividade através do desempenho de várias funções em um campo de trabalho, mas de estender ao ensino médio processos de trabalho reais, possibilitando -se a assimilação não apenas teórica, mas também prática, dos princípios científicos que estão na base da produção moderna ”.
Sendo assim, a formação desse \textbf{trabalhador} deve ser abrangente de forma que ele seja preparado para o mundo do trabalho e para a vida.
A definição do \textbf{perfil} do \textbf{profissional} estabelecido em alguns \textbf{cursos} \textbf{técnicos} permitirá a vivência de estágio ou aprendizagem \textbf{profissional}, visitas \textbf{técnicas}, como formas de dialogar com o mundo do trabalho.
Havendo estágio supervisionado, este deve ser concebido como uma prática educativa e como atividade curricular intencionalmente planejada, integrando o \textbf{currículo} do \textbf{curso} com \textbf{carga} \textbf{horária} acrescida ao mínimo estabelecido legalmente para a \textbf{habilitação} \textbf{profissional}.
Ao indicar caminhos para o desenvolvimento de competências e ampliar as perspectivas de formação permanente e ao longo da vida, a convergência curricular tem como pressuposto o aproveitamento de estudos e a valorização de experiências e aprendizagens já desenvolvidas, o que possibilita o aprimoramento \textbf{profissional} e incentiva o progressivo avanço dos níveis de \textbf{escolaridade}.
Dessa forma, o aluno pode imprimir ritmo e direção ao seu percurso \textbf{formativo}, com autonomia para acelerar, postergar, interromper ou retomar seus estudos, de acordo com as suas necessidades e as novas exigências \textbf{profissionais}.
Em síntese, a arquitetura curricular convergente torna real a \textbf{oferta} de Unidades Curriculares ou de \textbf{certificações} intermediárias que podem ser compartilhadas entre diferentes \textbf{cursos}, porque, em essência, são partes integrantes dos próprios \textbf{cursos}.
Tomadas como princípios para a elaboração dos desenhos curriculares, a flexibilidade e a convergência curricular ampliam as possibilidades \textbf{formativas} e favorecem a contextualização e atualização do \textbf{currículo}, já que essa concepção permite aos alunos construir uma \textbf{trajetória} \textbf{profissional} particular sem perder de vista as \textbf{diretrizes} dos \textbf{cursos} de formação, além de colaborar para que os egressos tenham uma formação sólida e abrangente.
Portanto, é a configuração de \textbf{itinerários} \textbf{formativos} pautada nesses dois princípios que torna factível aos alunos e às empresas a escolha entre múltiplas \textbf{trajetórias} \textbf{formativas} e a busca por formatos de \textbf{cursos} que se adaptem aos seus interesses e possibilidades.
É necessário também estar atento às mudanças tecnológicas, uma vez que essas alterações afetam a relação do homem com o próprio trabalho e com seu projeto de vida.
Diante do surgimento de novas ocupações, das mudanças nas formas de gestão das organizações e no próprio ambiente de trabalho, as pessoas buscam propostas \textbf{formativas} que permitam tanto complementar e desenvolver competências no âmbito das ocupações em que atuam como ascender a níveis mais elevados de formação \textbf{profissional}, com possibilidades de mudança de \textbf{carreira}.
A análise dos impactos desses fatores é fundamental para que a educação \textbf{profissional} seja uma grande aliada na construção do projeto de vida das pessoas.
Ao reconhecer o aluno como sujeito autônomo, portador de história e de saberes, a SEED/AP viabiliza a construção de \textbf{itinerários} \textbf{formativos} \textbf{alinhados} às expectativas dos indivíduos, oferecendo chances reais de inserção qualificada no mundo do trabalho.
Em síntese, as expectativas, as condições de laboralidade e as oportunidades de \textbf{carreira} devem orientar o desenvolvimento de \textbf{ofertas} cada vez mais diversificadas, o que exige a avaliação permanente das ações do Seed/AP.
Nesse contexto, a Pesquisa de Avaliação do Egresso e a Pesquisa Nacional de Qualidade Percebida dos \textbf{Cursos} da Seed/AP são ferramentas importantes para mapear os fatores que influenciam as prioridades e escolhas dos alunos.
Esse olhar criterioso do ponto de vista \textbf{técnico} e pedagógico é fundamental para a identificação de competências comuns entre as ocupações no decorrer da elaboração de Planos de \textbf{Curso} Nacionais, de forma a construir \textbf{trajetórias} de formação por ocupação e por segmento, com interface entre os \textbf{cursos} tanto na direção horizontal ( \textbf{cursos} do mesmo nível ) como na vertical ( \textbf{cursos} de níveis distintos ).
Daí a necessidade de os \textbf{currículos} serem permeáveis, de forma a \textbf{viabilizar} a convergência entre as Unidades Curriculares de um conjunto de formações.
Contudo, cabe salientar que, apesar de a ênfase da convergência curricular estar na competência, é preciso delinear um desenho curricular em que toda a organização interna do \textbf{curso} permita o aproveitamento contínuo de estudos ; ou seja, é preciso articular de forma dinâmica as competências, os elementos da competência, as \textbf{cargas} \textbf{horárias}, as estratégias metodológicas, os critérios de avaliação e os recursos didáticos, equipamentos e materiais.
Considerando o exposto, em termos de portfólio, os \textbf{Itinerários} \textbf{Formativos} Nacionais propõem um modelo de organização dos \textbf{cursos} e programas que permite conferir maior coerência, aderência e organicidade à \textbf{oferta} de educação \textbf{profissional} realizada pela Seed/AP.
Ao projetar \textbf{cursos} e programas educacionais orientados por \textbf{itinerários} \textbf{formativos}, a Seed/AP prioriza os percursos de profissionalização das pessoas, adotando metodologias educacionais que permitam superar a segmentação da organização curricular e contribuir para o contínuo e articulado aproveitamento de estudos e experiências \textbf{profissionais}.
Dessa forma, compreender a estrutura e a organização dos \textbf{cursos} é fundamental para planejar, classificar e orientar a \textbf{oferta} de educação \textbf{profissional} de forma coerente com a complexidade das ocupações identificadas nos \textbf{itinerários} \textbf{profissionais}, sem perder de vista os objetivos educacionais de cada programação.
Sendo a colaboração e a autonomia \textbf{prerrogativas} do Modelo Pedagógico Seed/AP, a implementação de \textbf{Itinerários} \textbf{Formativos} pelos Departamentos Regionais deve considerar as demandas locais e as orientações político-institucionais.
Em outras palavras, o Itinerário \textbf{Formativo} é um instrumento que visa incrementar e fortalecer a educação \textbf{profissional} oferecida pela SEED/AP, estabelecendo parâmetros para a organização do portfólio de cada Departamento Regional.

% matched lemmas: alinhar, atender, carga, carreira, certificação, currículo, curso, diretriz, escolaridade, formativo, habilitação, horário, itinerário, legislação, oferta, perfil, prerrogativa, profissional, trabalhador, trajetória, técnico, viabilizar
\end{textsample}
